\documentclass[10pt,a4paper]{amsart}
\usepackage[margin=19mm]{geometry}
\usepackage{amsmath,amssymb,mathtools}
\usepackage[scr=rsfs,cal=euler]{mathalfa}
\usepackage{xfrac}
\usepackage[unicode,hidelinks,bookmarks=false]{hyperref}
\newcommand{\ie}{i.e.\ }
\newcommand{\cf}{cf.\ }
\newcommand{\resp}{respectively\ }
\newcommand{\Subsection}{\subsection}
\providecommand{\nobreakdash}{\nobreak\hbox{-}}
\setlength{\emergencystretch}{2em}
\multlinegap0pt
\usepackage{bm}
\usepackage{cleveref}

\theoremstyle{plain}
\newtheorem{theorem}{Theorem}
\newtheorem{proposition}[theorem]{Proposition}
\newtheorem{lemma}[theorem]{Lemma}
\newtheorem{corollary}[theorem]{Corollary}
\theoremstyle{definition}
\newtheorem{definition}[theorem]{Definition}
\newtheorem{assumption}[theorem]{Assumption}
\theoremstyle{remark}
\newtheorem{remark}[theorem]{Remark}
\title[Safeguarded Accelerated Adaptive Search under Unreliable Gradient Oracles]{Safeguarded Accelerated Adaptive Search under Unreliable Gradient Oracles: Theorem 1}
\author{Shunzhi Zhang, Shichen Liao, Congying Han, Tiande Guo}
\date{Source: arXiv:2604.15526v2 (CC BY 4.0)}
\begin{document}
\maketitle

\noindent\emph{Source and licence.} The delivered work is the article shown in the title above, version arXiv:2604.15526v2 (CC BY 4.0), distributed under the Creative Commons Attribution 4.0 International licence, as recorded in the arXiv licence metadata for this exact version and shown on the abstract page saved with this item. The full licence notice, which carries the licence URL and the address of that abstract page, is the bundled file \texttt{licenses/arxiv-2604.15526-CC-BY-4.0.txt}.\par\smallskip

\noindent\textit{Editorial scope.} Counted unit: the article's Theorem 1 (source0.tex lines 1263--1331, source label \texttt{Big\_threorem}); statement and proof are the article's own text line for line, in the article's own notation ($\Phi_t$, $T_{\boldsymbol\varepsilon}$, $H(t)$, $P(t)$, $\varepsilon_{\mathrm{res}}$), with no line omitted, substituted or re-flowed. Required context retained uncounted: the article's stopping-time definition and the surrounding setup (source0.tex lines 642--669), which the counted statement refers to through \texttt{\textbackslash Cref}. Editorial formatting only: an AMS \texttt{amsart} preamble declaring the theorem environments the retained lines use, with \texttt{bm} loaded for the article's bold math and \texttt{cleveref} loaded so that the article's \texttt{\textbackslash Cref} cross-reference resolves; the article ships no class file of its own and the corpus does not redistribute it, so none is shipped or input; sequential numbering of the retained environments in order of appearance, whereas the article prints the counted theorem as Theorem 1; equation numbering is not reproduced, and every cross-reference inside the retained text resolves to the wrapper's numbering. No citation occurs inside any retained line, so the wrapper carries no bibliography. No asset-inclusion command, no drawing environment and no image asset occurs in this package.


\section*{Retained context: Stopping\_time (source0.tex lines 642--669)}

\begin{definition}[\(\boldsymbol{\varepsilon}\)-stopping time]
\label{Stopping_time}
For an  accuracy target
\(\boldsymbol{\varepsilon}
=(\varepsilon_\phi,\varepsilon_\nabla)\), where
\(\varepsilon_\phi>0\) and \(\varepsilon_\nabla\ge0\), define
\begin{equation}\label{eq_dfn_stopping_time}
T_{\boldsymbol{\varepsilon}}
\coloneqq 
\inf\left\{
t\ge1:
\min\{\phi(\hat {\bm{x}}_{t+1}),\phi(\hat {\bm{y}}_t)\}-\phi^\star
\le\varepsilon_\phi
\ \text{or}\
\|\nabla\phi(\hat {\bm{y}}_t)\|\le\varepsilon_\nabla
\right\}.
\end{equation}
We use the convention $\inf\varnothing=+\infty$. The definition includes both the extrapolated point and the candidate, whether or not the trial is accepted. The event that it holds at trial $t$ is $\mathcal F_{t-1/2}$-measurable, so $T_{\boldsymbol{\varepsilon}}$ is a stopping time with respect to $\{\mathcal F_t\}_{t\geq0}$. When $\varepsilon_\nabla=0$, we abbreviate $\varepsilon_\phi$ by $\varepsilon$ and write
\[
T_\varepsilon
\coloneqq 
\inf\left\{
t\ge1:
\min\{\phi(\hat {\bm{x}}_{t+1}),\phi(\hat {\bm{y}}_t)\}-\phi^\star
\le\varepsilon
\right\}.
\]
\end{definition}


\section{The counted unit: Theorem 1 and its complete original proof}

\begin{theorem}[A high-probability stopping-time bound]
\label{Big_threorem}
Let
\(\boldsymbol\varepsilon=(\varepsilon_\phi,\varepsilon_\nabla)\), where
\(\varepsilon_\phi>0\) and \(\varepsilon_\nabla\ge0\), and let
\(T_{\boldsymbol\varepsilon}\) be the stopping time in
\Cref{Stopping_time}.  Suppose that \(\Phi_t\ge0\) is adapted and
compatible with the stopping rule in the sense that
$\Phi_t\le\varepsilon_\phi
\Longrightarrow
T_{\boldsymbol\varepsilon}\le t$.
Assume that \(\Phi_1\le\Phi_0\) and that there exist adapted sequences
\(V_t\ge0\) and \(R_t\in\mathbb R\) satisfying
\begin{equation}\label{eq:abstract_one_step}
\Phi_{t+1}\le V_t\Phi_t+R_t,
\qquad 1\le t<T_{\boldsymbol\varepsilon}.
\end{equation}
Suppose further that there exist
\(H\colon \mathbb N^+\to\mathbb R_+\),
\(P\colon \mathbb N^+\to[0,1]\),
\(\varepsilon_{\mathrm{res}}\ge0\), and
\(\bar T\in\mathbb N^+\) such that, for every integer
\(t\ge\bar T\),
\begin{equation}\label{eq:abstract_tail_event}
\begin{aligned}
\mathbb P\Bigl(&
\Bigl\{
 t<T_{\boldsymbol\varepsilon},\
 \prod_{i=1}^tV_i>H(t)
\Bigr\}\cup\,
\Bigl\{
 t<T_{\boldsymbol\varepsilon},\
 \sum_{i=1}^t
 \bigl(\prod_{j=i+1}^tV_j\bigr)R_i
 >\varepsilon_{\mathrm{res}}
\Bigr\}
\Bigr)
\le P(t).
\end{aligned}
\end{equation}
Assume that \(H\) is non-increasing on
\(\{t\in\mathbb N^+:t\ge\bar T\}\) and that
\(\varepsilon_\phi>\varepsilon_{\mathrm{res}}\). Define its generalized
inverse by\footnote{We use the conventions that an empty product equals
one and \(\inf\varnothing=+\infty\).}
\begin{equation}\label{eq:abstract_inverse_H}
H^{-1}(u)
\coloneqq 
\inf\left\{
s\in\mathbb N^+:
s\ge\bar T,\ H(s)\le u
\right\},
\qquad u>0.
\end{equation}
For \(\Phi_0>0\),\footnote{If
\(\Phi_0=0\), then \(\Phi_1=0\) and the conclusion follows immediately
from compatibility.} every integer \(t\) satisfying
\begin{equation}\label{eq:abstract_complexity_inverse}
t\ge\max\left\{
\bar T,\,
H^{-1}\!\left(
\frac{\varepsilon_\phi-\varepsilon_{\mathrm{res}}}{\Phi_0}
\right)
\right\}
\end{equation}
guarantees
\(\mathbb P\{T_{\boldsymbol\varepsilon}\leq t+1\}\geq1-P(t)\).

\end{theorem}

\begin{proof}
\phantomsection\label{proof_bigtheorem}
On $\{t<T_{\boldsymbol\varepsilon}\}$, iterate \eqref{eq:abstract_one_step}, using $V_i\geq0$ and $\Phi_1\leq\Phi_0$, to obtain
\begin{equation}\label{eq:abstract_lyapunov_iteration}
\textstyle\Phi_{t+1}
\le
\Phi_0\prod_{i=1}^tV_i
+
\sum_{i=1}^t
\left(\prod_{j=i+1}^tV_j\right)R_i.
\end{equation}
If \(t\) satisfies \eqref{eq:abstract_complexity_inverse}, the
definition of \(H^{-1}\) and the monotonicity of \(H\) give
\(H(t)\Phi_0+\varepsilon_{\mathrm{res}}\leq\varepsilon_\phi\).
On \(\{t<T_{\boldsymbol\varepsilon}\}\) outside the event in
\eqref{eq:abstract_tail_event},
\eqref{eq:abstract_lyapunov_iteration} therefore gives
\(\Phi_{t+1}\leq\varepsilon_\phi\), and compatibility yields
\(T_{\boldsymbol\varepsilon}\leq t+1\).
The same conclusion holds on \(\{T_{\boldsymbol\varepsilon}\leq t\}\).
Hence \(\mathbb P\{T_{\boldsymbol\varepsilon}\leq t+1\}\geq1-P(t)\).
\end{proof}

\end{document}
